% OWNER: sections/s0_cover.tex. Edit only this file.
% Cover: same coordinate grid as the 7709 report cover.tex (mm from page top-left, text baselines).
% Title Arial Bold 14pt, subtitle 10pt grey; right column Quant Team + authors (name 8pt bold, e-mail +3.9mm, pitch 9.5mm; no phone numbers in the public version).
% \covertitlelines (1 or 2) must match the typeset title; the subtitle moves down 5.98mm per extra title line.
\thispagestyle{dartpage}
\begin{tikzpicture}[remember picture,overlay,x=1mm,y=-1mm,shift={(current page.north west)}]
  \node[anchor=base west,inner sep=0,text width=136.2mm,align=left] at (14.54,32.69)
    {\fontsize{14}{17}\selectfont\bfseries\addfontfeatures{LetterSpace=0.64}\color{navy}\RaggedRight\hyphenpenalty=10000 \exhyphenpenalty=10000 \relax\covertitle};
  \node[anchor=base west,inner sep=0,text width=136.2mm,align=left] at (14.54,{40.01+5.98*(\covertitlelines-1)})
    {\fontsize{10}{12}\selectfont\addfontfeatures{LetterSpace=0.41}\color{subgrey}\coversubtitle};
  \node[anchor=base west,inner sep=0] at (154.57,31.55) {\kofont\bfseries\fontsize{10}{10}\selectfont\addfontfeatures{LetterSpace=4.2}\color{navy}Quant Team};
  \fill[navy] (154.13,33.15) rectangle (197.73,33.49);
  \def\analystlist{38.02/Choi Munseok/chlanstjr0301@gmail.com,
                   47.52/Choi Yeayun/yunn1215@gmail.com}
  \foreach \y/\n/\e in \analystlist {
    \node[anchor=base west,inner sep=0] at (154.57,\y) {\kofont\bfseries\fontsize{8}{8}\selectfont\addfontfeatures{LetterSpace=3.0}\color{black}\n};
    \node[anchor=base west,inner sep=0] at (154.57,\y+3.9) {\kofont\fontsize{7}{7}\selectfont\addfontfeatures{LetterSpace=2.1}\color{body}\e};
  }
\end{tikzpicture}%
\vspace*{\dimexpr 5.98mm*(\covertitlelines-1)+24.8mm\relax}%
\par
\newcommand{\coverblock}[2]{{\fontsize{9.5}{11.6}\selectfont\bfseries\color{navy}\RaggedRight #1\par}\vspace{0.5pt}{\fontsize{9}{11}\selectfont #2\par}\vspace{2.5pt}}
\coverblock{This design document ties every wrapper decision to a measurement or a stated rule}{This is the software design document of \code{pytimefolio}: it records what the competition server does (Section~\ref{sec:layers}), the design that follows from it (Section~\ref{sec:design}) and the live tests (Section~\ref{sec:live}); \figref{summary} maps each layer to its decisions. Timefolio publishes no API for its Road to Fund Manager competition, so the wrapper sends the requests the competition's web client sends; we reproduced those requests and measured the server's replies on 24 Nov 2026. The organiser stated by e-mail that it does not prohibit automated orders and that an IP address that places 100 or more orders in a day may be blocked [5], so the wrapper stops at 80.}
\coverblock{Measurements set the design: one reused connection, one token and a 3\,s poll}{Section~\ref{sec:layers} measures 5 protocol layers. A REST request takes 16.0\ms{} on a reused HTTP/2 connection and 106.5\ms{} on a new one. The access token lasts 900\,s; the cookie jar set at login lasts 7 days, and neither renewal nor logging out changes it, so the wrapper guards the jar as a credential. The live feed sends nothing unless asked, so the wrapper asks every 3\,s. A request that echoes the previous reply's stamp (\code{t}) exactly gets a delta of the changed fields (1,053 to 1,912 bytes in the session, against an 11,344-byte snapshot); any other request gets a full snapshot.}
\coverblock{No order ID comes back, so an order the server may have received is never resent}{Orders (Section~\ref{sec:orders}) are placed over REST and confirmed in the feed, and Section~\ref{sec:design} builds the wrapper around this safety rule. The wrapper resends an order once only when the connection was never established, or after a 401, which ASP.NET returns before the order handler runs; otherwise it looks for the order in the feed for 15\,s, then in the order list, then raises an error so that the caller can check the web page. Before every order it checks the feed's age (under 26\,s), the business day and the stock limit; \code{plan()} checks the target weights' sector and total limits.}
\coverblock{Live tests exposed 2 defects, since fixed, and showed that share counts round down}{In Section~\ref{sec:live}'s live tests, start-up, the feed, a buy and a sell all worked; 31 offline tests cover failures a live server cannot be made to produce. The live tests exposed 2 defects: \code{plan()} accepted targets summing to 139\%, and \code{guidelines()} looked up sector limits by the wrong key. An order of 0.04\%, worth 1.76 shares, was sent twice; both times the server set the shares ordered (\code{qty2ord}) to 1, which rules out rounding to the nearest share, so \code{plan()} drops orders below one share.}
\vspace{2pt}

\figinlinetrue% OWNER: figs/summary.tex. Edit only this file.
\fig[wide]{summary}{Each layer was measured live, and each finding fixes the decisions listed in Appendix~B}[Main finding per section and the decisions it supports (D-numbers of Appendix~\ref{app:decisions})]{%
{\tabfont\begin{tabularx}{\linewidth}{>{\RaggedRight\arraybackslash}p{27mm} >{\RaggedRight\arraybackslash}X >{\RaggedRight\arraybackslash}p{29mm}}
\thd{Layer} & \thd{Finding} & \thd{Decisions}\\
\textbf{\ref{sec:transport} Transport} & A request takes 106.5\ms{} on a new HTTP/2 connection and 16.0\ms{} on a reused one; the server closes an idle connection after 15 to 30\,s & D1, D2, D18\\\rowsep
\textbf{\ref{sec:auth} Login session} & The access token lasts 900\,s and is accepted 300\,s past expiry; the cookie jar lasts 7 days, renewal never extends it and \code{Logout} revokes nothing & D3 to D5, D16, D17, D26, D35, D36\\\rowsep
\textbf{\ref{sec:queries} Queries} & A Bearer header authenticates every query; a non-trading date returns another day's data with HTTP~200 & D6, D7, D27\\\rowsep
\textbf{\ref{sec:feed} Live feed} & The feed answers only when asked (every 3\,s); an exact \code{t} echo gets a 1,053 to 1,912-byte delta instead of the 11,344-byte snapshot; an open WebSocket outlives its token & D8, D9, D14, D30, D31\\\rowsep
\textbf{\ref{sec:orders} Orders} & \code{AddOrder} returns no order ID; \code{wei} is the order size (0.3\% asked for 7 shares); every finished order ends in \code{Done} & D11, D13, D15, D29\\\rowsep
\textbf{\ref{sec:design} Wrapper design} & Direct HTTP with asyncio keeps one login session; strategy and market data stay with the caller; the organiser may block an IP address at 100 orders a day & D10, D12, D33, D34\\\rowsep
\textbf{\ref{sec:live} Live tests} & Start-up, feed, a buy and a sell worked; \code{qty2ord} rounds down (1.76 shares' worth gave 1); \code{plan()} needed a 100\% cap and a sector check & D19 to D25, D28, D32\\
\end{tabularx}}}{Measured 05:38 to 10:34 KST on 24 Nov 2026 with one account; the market opened at 09:00. Timings in Section~\ref{sec:transport} were taken before the open.}{DART Quant Team}
\figinlinefalse

\clearpage
